% Generated draft: requires mathematical and source review.
\documentclass[11pt]{article}
\usepackage[margin=25mm]{geometry}
\usepackage{fontspec}
\setmainfont{Latin Modern Roman}
\usepackage{amsmath,amssymb,mathtools,unicode-math}
\setmathfont{Latin Modern Math}
\usepackage{enumitem,needspace,xurl,hyperref}
\setlength{\parindent}{0pt}
\setlength{\parskip}{5pt}
\setlength{\emergencystretch}{3em}
\newcommand{\N}{\mathbb{N}}
\newcommand{\Z}{\mathbb{Z}}
\newcommand{\Q}{\mathbb{Q}}
\newcommand{\R}{\mathbb{R}}
\newcommand{\C}{\mathbb{C}}
\newcommand{\F}{\mathbb{F}}
\newcommand{\A}{\mathbb{A}}
\newcommand{\PP}{\mathbb{P}}
\newcommand{\E}{\mathbb{E}}
\newcommand{\eps}{\varepsilon}
\newcommand{\dd}{\,\mathrm{d}}
\begin{document}
\section*{UnsolvedMath ID: 20000745}
\textbf{Source problem code:} catalogue entry
\subsection{Definitions and mathematical statement}
Let \(p_1=2<p_2<\cdots\) be the consecutive primes, \(d_n=p_{n+1}-p_n\), and
\[
 A_n=p_{n+2}-2p_{n+1}+p_n=d_{n+1}-d_n.
\]
T. D. Wooley's original question asks how often \(A_n\) is small or large, without specifying a threshold or an asymptotic frequency target. For example the quantities to study may be \(\#\{n:p_n\leq X,\ |A_n|\leq Y\}\) and separate counts of \(A_n\geq B(p_n)\) or \(A_n\leq-B(p_n)\), but no particular \(Y\) or \(B\) is imposed by the source. It reports that Erd\H{o}s proved \(|A_n|\geq(1+\delta)\log p_n\) infinitely often for some \(\delta>0\), and records A. Balog's informal expectation that the magnitude can match the then-largest known prime-gap scale. The latter statement is not supplied as a precise theorem.

For the aid's explicit large-value assertions, let \(\log_1x=\log x\), \(\log_{j+1}x=\log(\log_jx)\), and work with sufficiently large \(x\) that all displayed iterates are positive. Put
\[
 L(x)=\frac{\log x\,\log_2x\,\log_4x}{\log_3x},\qquad
 F(x)=\frac{\log x\,\log_2x\,\log_5x}{\log_3x}.
\]
The supplied Sun--Pan corollary is that there exist \(c,c_0>0\) and \(X_0\) such that, for every \(X\geq X_0\),
\[
 \#\{n:p_n\leq X,\ A_n\geq cL(p_n)\}\geq c_0\log X.
\]
The report emphasizes retaining the every-sufficiently-large-scale uniformity of the underlying bounded-cluster/long-gap construction; an assertion only along an unspecified subsequence would not give this count as stated.

A separate reported specialization of Baker--Freiberg's two-gap chain theorem, Theorem 6.4, using scaled convex and concave configurations of 64 points, yields
\[
 \limsup_{n\to\infty}\frac{A_n}{F(p_n)}=+\infty,
 \qquad
 \liminf_{n\to\infty}\frac{A_n}{F(p_n)}=-\infty.
\]
The report includes a denominator-transfer check to ensure that the comparison scale is evaluated at \(p_n\), rather than an auxiliary construction parameter. Its configuration coordinates, theorem hypotheses, and constants are not reproduced in the supplied record.

The aid also attributes to Pintz counts of order at least \(\log X\) for occurrences of each sign on a ``classical Erd\H{o}s--Rankin scale.'' That scale and its quantitative constants are not defined here, so this attribution is preserved without substituting an invented logarithmic profile. Finally it attributes to Banks--Freiberg--Turnage-Butterbaugh a fixed-tuple argument giving some fixed integers \(a_+>0\), \(a_-<0\), each occurring as \(A_n\) infinitely often. These values are necessarily even once all relevant primes are odd. No claim is made that every nonzero even value recurs, that \(A_n=0\) recurs, or that any such recurrence has a specified positive density. These occurrence results are partial information toward the broad small/large frequency problem.
\subsection{Short English statement}
The source asks about the sizes and frequencies of consecutive-prime second differences. The supplied partial results give logarithmically many positive values on an explicit large-gap scale, both signed unbounded normalized extremes on another explicit scale, and recurrence of some fixed nonzero value of each sign. They do not give a complete distribution or small-value classification.
\subsection{Sources}
\begin{itemize}
\item Local registry: permanent UnsolvedMath ID \(20000745\), exact public code \texttt{catalogue entry}; \texttt{aim-analytic-number-theory-notes.json}, \texttt{notes[108]}.
\item AIM workshop notes \emph{Gaps Between Primes}, December 2005, recorded by Angel Kumchev, Problem 9, ``Second order differences between primes'': \url{https://aimath.org/WWN/primegaps/WorkshopProblems.pdf}. This URL is supplied provenance; no PDF access is claimed.
\item Original attributions: T. D. Wooley for the question, Erd\H{o}s for inequality (3), and A. Balog for the informal largest-gap-scale comment. No complete paper citations are supplied.
\item Supplied research summary, ``Signed large and recurrent fixed second differences of consecutive primes,'' attempt 1, invokes Sun--Pan, Baker--Freiberg Theorem 6.4, Pintz, and Banks--Freiberg--Turnage-Butterbaugh without paper titles, publication details, or complete theorem statements. Original remarks and literature fields are empty.
\end{itemize}
\end{document}
